\IfFileExists{stacks-project.cls}{%
\documentclass{stacks-project}
}{%
\documentclass{amsart}
}

% For dealing with references we use the comment environment
\usepackage{verbatim}
\newenvironment{reference}{\comment}{\endcomment}
%\newenvironment{reference}{}{}
\newenvironment{slogan}{\comment}{\endcomment}
\newenvironment{history}{\comment}{\endcomment}

% For commutative diagrams we use Xy-pic
\usepackage[all]{xy}

% We use 2cell for 2-commutative diagrams.
\xyoption{2cell}
\UseAllTwocells

% We use multicol for the list of chapters between chapters
\usepackage{multicol}

% This is generally recommended for better output
\usepackage{lmodern}
\usepackage[T1]{fontenc}

% For cross-file-references
\usepackage{xr-hyper}

% Package for hypertext links:
\usepackage{hyperref}

% For any local file, say "hello.tex" you want to link to please
% use \externaldocument[hello-]{hello}
\externaldocument[introduction-]{introduction}
\externaldocument[conventions-]{conventions}
\externaldocument[sets-]{sets}
\externaldocument[categories-]{categories}
\externaldocument[topology-]{topology}
\externaldocument[sheaves-]{sheaves}
\externaldocument[sites-]{sites}
\externaldocument[stacks-]{stacks}
\externaldocument[fields-]{fields}
\externaldocument[algebra-]{algebra}
\externaldocument[brauer-]{brauer}
\externaldocument[homology-]{homology}
\externaldocument[derived-]{derived}
\externaldocument[simplicial-]{simplicial}
\externaldocument[more-algebra-]{more-algebra}
\externaldocument[smoothing-]{smoothing}
\externaldocument[modules-]{modules}
\externaldocument[sites-modules-]{sites-modules}
\externaldocument[injectives-]{injectives}
\externaldocument[cohomology-]{cohomology}
\externaldocument[sites-cohomology-]{sites-cohomology}
\externaldocument[dga-]{dga}
\externaldocument[dpa-]{dpa}
\externaldocument[sdga-]{sdga}
\externaldocument[hypercovering-]{hypercovering}
\externaldocument[schemes-]{schemes}
\externaldocument[constructions-]{constructions}
\externaldocument[properties-]{properties}
\externaldocument[morphisms-]{morphisms}
\externaldocument[coherent-]{coherent}
\externaldocument[divisors-]{divisors}
\externaldocument[limits-]{limits}
\externaldocument[varieties-]{varieties}
\externaldocument[topologies-]{topologies}
\externaldocument[descent-]{descent}
\externaldocument[perfect-]{perfect}
\externaldocument[more-morphisms-]{more-morphisms}
\externaldocument[flat-]{flat}
\externaldocument[groupoids-]{groupoids}
\externaldocument[more-groupoids-]{more-groupoids}
\externaldocument[etale-]{etale}
\externaldocument[chow-]{chow}
\externaldocument[intersection-]{intersection}
\externaldocument[pic-]{pic}
\externaldocument[weil-]{weil}
\externaldocument[adequate-]{adequate}
\externaldocument[dualizing-]{dualizing}
\externaldocument[duality-]{duality}
\externaldocument[discriminant-]{discriminant}
\externaldocument[derham-]{derham}
\externaldocument[local-cohomology-]{local-cohomology}
\externaldocument[algebraization-]{algebraization}
\externaldocument[curves-]{curves}
\externaldocument[resolve-]{resolve}
\externaldocument[models-]{models}
\externaldocument[functors-]{functors}
\externaldocument[equiv-]{equiv}
\externaldocument[pione-]{pione}
\externaldocument[etale-cohomology-]{etale-cohomology}
\externaldocument[proetale-]{proetale}
\externaldocument[relative-cycles-]{relative-cycles}
\externaldocument[more-etale-]{more-etale}
\externaldocument[trace-]{trace}
\externaldocument[crystalline-]{crystalline}
\externaldocument[spaces-]{spaces}
\externaldocument[spaces-properties-]{spaces-properties}
\externaldocument[spaces-morphisms-]{spaces-morphisms}
\externaldocument[decent-spaces-]{decent-spaces}
\externaldocument[spaces-cohomology-]{spaces-cohomology}
\externaldocument[spaces-limits-]{spaces-limits}
\externaldocument[spaces-divisors-]{spaces-divisors}
\externaldocument[spaces-over-fields-]{spaces-over-fields}
\externaldocument[spaces-topologies-]{spaces-topologies}
\externaldocument[spaces-descent-]{spaces-descent}
\externaldocument[spaces-perfect-]{spaces-perfect}
\externaldocument[spaces-more-morphisms-]{spaces-more-morphisms}
\externaldocument[spaces-flat-]{spaces-flat}
\externaldocument[spaces-groupoids-]{spaces-groupoids}
\externaldocument[spaces-more-groupoids-]{spaces-more-groupoids}
\externaldocument[bootstrap-]{bootstrap}
\externaldocument[spaces-pushouts-]{spaces-pushouts}
\externaldocument[spaces-chow-]{spaces-chow}
\externaldocument[groupoids-quotients-]{groupoids-quotients}
\externaldocument[spaces-more-cohomology-]{spaces-more-cohomology}
\externaldocument[spaces-simplicial-]{spaces-simplicial}
\externaldocument[spaces-duality-]{spaces-duality}
\externaldocument[formal-spaces-]{formal-spaces}
\externaldocument[restricted-]{restricted}
\externaldocument[spaces-resolve-]{spaces-resolve}
\externaldocument[formal-defos-]{formal-defos}
\externaldocument[defos-]{defos}
\externaldocument[cotangent-]{cotangent}
\externaldocument[examples-defos-]{examples-defos}
\externaldocument[algebraic-]{algebraic}
\externaldocument[examples-stacks-]{examples-stacks}
\externaldocument[stacks-sheaves-]{stacks-sheaves}
\externaldocument[criteria-]{criteria}
\externaldocument[artin-]{artin}
\externaldocument[quot-]{quot}
\externaldocument[stacks-properties-]{stacks-properties}
\externaldocument[stacks-morphisms-]{stacks-morphisms}
\externaldocument[stacks-limits-]{stacks-limits}
\externaldocument[stacks-cohomology-]{stacks-cohomology}
\externaldocument[stacks-perfect-]{stacks-perfect}
\externaldocument[stacks-introduction-]{stacks-introduction}
\externaldocument[stacks-more-morphisms-]{stacks-more-morphisms}
\externaldocument[stacks-geometry-]{stacks-geometry}
\externaldocument[moduli-]{moduli}
\externaldocument[moduli-curves-]{moduli-curves}
\externaldocument[examples-]{examples}
\externaldocument[exercises-]{exercises}
\externaldocument[guide-]{guide}
\externaldocument[desirables-]{desirables}
\externaldocument[coding-]{coding}
\externaldocument[obsolete-]{obsolete}
\externaldocument[fdl-]{fdl}
\externaldocument[index-]{index}

% Theorem environments.
%
\theoremstyle{plain}
\newtheorem{theorem}[subsection]{Theorem}
\newtheorem{proposition}[subsection]{Proposition}
\newtheorem{lemma}[subsection]{Lemma}

\theoremstyle{definition}
\newtheorem{definition}[subsection]{Definition}
\newtheorem{example}[subsection]{Example}
\newtheorem{exercise}[subsection]{Exercise}
\newtheorem{situation}[subsection]{Situation}

\theoremstyle{remark}
\newtheorem{remark}[subsection]{Remark}
\newtheorem{remarks}[subsection]{Remarks}

\numberwithin{equation}{subsection}

% Macros
%
\def\lim{\mathop{\mathrm{lim}}\nolimits}
\def\colim{\mathop{\mathrm{colim}}\nolimits}
\def\Spec{\mathop{\mathrm{Spec}}}
\def\Hom{\mathop{\mathrm{Hom}}\nolimits}
\def\Ext{\mathop{\mathrm{Ext}}\nolimits}
\def\SheafHom{\mathop{\mathcal{H}\!\mathit{om}}\nolimits}
\def\SheafExt{\mathop{\mathcal{E}\!\mathit{xt}}\nolimits}
\def\Sch{\mathit{Sch}}
\def\Mor{\mathop{\mathrm{Mor}}\nolimits}
\def\Ob{\mathop{\mathrm{Ob}}\nolimits}
\def\Sh{\mathop{\mathit{Sh}}\nolimits}
\def\NL{\mathop{N\!L}\nolimits}
\def\CH{\mathop{\mathrm{CH}}\nolimits}
\def\proetale{{pro\text{-}\acute{e}tale}}
\def\etale{{\acute{e}tale}}
\def\QCoh{\mathit{QCoh}}
\def\Ker{\mathop{\mathrm{Ker}}}
\def\Im{\mathop{\mathrm{Im}}}
\def\Coker{\mathop{\mathrm{Coker}}}
\def\Coim{\mathop{\mathrm{Coim}}}

% Boxtimes
%
\DeclareMathSymbol{\boxtimes}{\mathbin}{AMSa}{"02}

%
% Macros for moduli stacks/spaces
%
\def\QCohstack{\mathcal{QC}\!\mathit{oh}}
\def\Cohstack{\mathcal{C}\!\mathit{oh}}
\def\Spacesstack{\mathcal{S}\!\mathit{paces}}
\def\Quotfunctor{\mathrm{Quot}}
\def\Hilbfunctor{\mathrm{Hilb}}
\def\Curvesstack{\mathcal{C}\!\mathit{urves}}
\def\Polarizedstack{\mathcal{P}\!\mathit{olarized}}
\def\Complexesstack{\mathcal{C}\!\mathit{omplexes}}
% \Pic is the operator that assigns to X its picard group, usage \Pic(X)
% \Picardstack_{X/B} denotes the Picard stack of X over B
% \Picardfunctor_{X/B} denotes the Picard functor of X over B
\def\Pic{\mathop{\mathrm{Pic}}\nolimits}
\def\Picardstack{\mathcal{P}\!\mathit{ic}}
\def\Picardfunctor{\mathrm{Pic}}
\def\Deformationcategory{\mathcal{D}\!\mathit{ef}}
\usepackage{cleveref}
\crefname{originallemma}{original lemma}{original lemmas}
\Crefname{originallemma}{Original lemma}{Original lemmas}
\makeatletter
\newlabel{lemma-faithful}{{0240}{source}{}{sourcepointer-1}{}}
\newlabel{lemma-faithful@cref}{{[originallemma][0][]0240}{[1][0][]source}{}{}{}}
\makeatother
\crefname{originallemma}{original lemma}{original lemmas}
\Crefname{originallemma}{Original lemma}{Original lemmas}
\makeatletter
\newlabel{lemma-ff-base-change-faithful}{{02VX}{source}{}{sourcepointer-2}{}}
\newlabel{lemma-ff-base-change-faithful@cref}{{[originallemma][0][]02VX}{[1][0][]source}{}{}{}}
\makeatother
\crefname{originallemma}{original lemma}{original lemmas}
\Crefname{originallemma}{Original lemma}{Original lemmas}
\makeatletter
\newlabel{lemma-fpqc-universal-effective-epimorphisms}{{023Q}{source}{}{sourcepointer-3}{}}
\newlabel{lemma-fpqc-universal-effective-epimorphisms@cref}{{[originallemma][0][]023Q}{[1][0][]source}{}{}{}}
\makeatother
\crefname{originallemma}{original lemma}{original lemmas}
\Crefname{originallemma}{Original lemma}{Original lemmas}
\makeatletter
\newlabel{lemma-pullback}{{023Y}{source}{}{sourcepointer-4}{}}
\newlabel{lemma-pullback@cref}{{[originallemma][0][]023Y}{[1][0][]source}{}{}{}}
\makeatother
\crefname{originalremark}{original remark}{original remarks}
\Crefname{originalremark}{Original remark}{Original remarks}
\makeatletter
\newlabel{remark-easier}{{02VP}{source}{}{sourcepointer-5}{}}
\newlabel{remark-easier@cref}{{[originalremark][0][]02VP}{[1][0][]source}{}{}{}}
\makeatother
\title{The pullback functor on descent data for schemes is fully faithful when the covering is fpqc over the same base}
\author{The Stacks Project Authors}
\date{stacks-project:descent.tex:local-snapshot-20261012; GNU FDL 1.2 or later}
\begin{document}
\maketitle
\noindent\textit{Editorial note.} All delivered mathematical text is the author's own, copied from the pinned Stacks Project chapter file descent.tex; no statement or proof step was rewritten, reordered or removed.\par\medskip
\noindent\textit{Editorial note.} The proof contains the author's sentence advising the reader to prove the displayed claim themselves rather than read the rest of the paragraph, and inviting an email if a cleaner argument is found. This is a reading suggestion, not an omission: the paragraph continues and supplies the pointwise calculation in full, so the card is collected and the sentence is reproduced verbatim. The step it refers to is recorded as a load-bearing obligation.\par\medskip
\noindent\textit{Editorial note.} The statement begins by referring to the situation of an earlier lemma of the chapter, which fixes the descent data and the pullback functor in question; that reference lies outside the counted span and is carried as an external reference stub rather than reproduced, and the definition it depends on is recorded as a standard prerequisite.\par\medskip
\noindent\textit{Editorial note.} The proof displays four commutative diagrams; all of them lie inside the counted span and are reproduced verbatim.\par\medskip
\noindent\textit{Editorial note.} The counted spans carry 6 navigation stubs; no bibliography citation occurs in the spans.\par\medskip
\noindent\textit{Editorial note.} The preamble span points at the chapter's own preamble.tex shipped with the frozen source, used verbatim as staged.\par\medskip
\noindent\textit{Allowed standard prerequisites (editorial).}\begin{itemize}
\item The situation of the pullback lemma: descent data relative to a morphism of schemes and the pullback functor between them. Same chapter, Lemma lemma-pullback, referred to in the statement at line 8779 and again at proof lines 8840-8841. Fixes the category of descent data and the functor whose full faithfulness is asserted, and supplies the description of the pulled back descent isomorphism.
\item Faithfulness of the pullback functor when the morphism is surjective and flat. Same chapter, Lemma lemma-faithful, invoked at proof lines 8790-8792. Gives faithfulness, reducing the lemma to the essential surjectivity on morphisms.
\item The pulled back descent data are base changes of the original data. Same chapter, Remark remark-easier, invoked at proof lines 8800-8803. Identifies the source of the pulled back datum with the base change of the source and the pulled back descent isomorphism with the base change of the original one.
\item Fpqc coverings are universal effective epimorphisms. Same chapter, Lemma lemma-fpqc-universal-effective-epimorphisms, invoked at proof lines 8855-8861. Factors the given morphism of descent data over X' through a unique morphism over X.
\item Base change along a surjective flat morphism is faithful. Same chapter, Lemma lemma-ff-base-change-faithful, invoked at proof lines 8871-8873. Lets the commutativity of the descent data square be checked after base change to the fibre product of X' with itself.
\end{itemize}\medskip
\ifdefined\hypertarget\hypertarget{sourcepointer-1}{}\fi
\noindent\textit{Source reference.} Original lemma 0240: sources/stacks-descent/pinned-original-source.tex, source lines 8758--8758. Pointer outside the counted span, invoked once for context; the load-bearing argument of the counted statement is the author's proof reproduced in full in the proof span.\par\medskip
\ifdefined\hypertarget\hypertarget{sourcepointer-2}{}\fi
\noindent\textit{Source reference.} Original lemma 02VX: sources/stacks-descent/pinned-original-source.tex, source lines 8724--8724. Pointer outside the counted span, invoked once for context; the load-bearing argument of the counted statement is the author's proof reproduced in full in the proof span.\par\medskip
\ifdefined\hypertarget\hypertarget{sourcepointer-3}{}\fi
\noindent\textit{Source reference.} Original lemma 023Q: sources/stacks-descent/pinned-original-source.tex, source lines 4348--4348. Pointer outside the counted span, invoked once for context; the load-bearing argument of the counted statement is the author's proof reproduced in full in the proof span.\par\medskip
\ifdefined\hypertarget\hypertarget{sourcepointer-4}{}\fi
\noindent\textit{Source reference.} Original lemma 023Y: sources/stacks-descent/pinned-original-source.tex, source lines 8496--8496. Pointer outside the counted span, invoked once for context; the load-bearing argument of the counted statement is the author's proof reproduced in full in the proof span.\par\medskip
\ifdefined\hypertarget\hypertarget{sourcepointer-5}{}\fi
\noindent\textit{Source reference.} Original remark 02VP: sources/stacks-descent/pinned-original-source.tex, source lines 8364--8364. Pointer outside the counted span, invoked once for context; the load-bearing argument of the counted statement is the author's proof reproduced in full in the proof span.\par\medskip
\begin{lemma}
\label{lemma-fully-faithful}
In the situation of Lemma \ref{lemma-pullback}
assume
\begin{enumerate}
\item $\{f : X' \to X\}$ is an fpqc covering (for example if $f$ is
surjective, flat, and quasi-compact), and
\item $S = S'$.
\end{enumerate}
Then the pullback functor is fully faithful.
\end{lemma}

\begin{proof}
Assumption (1) implies that $f$ is surjective and flat.
Hence the pullback functor is faithful by
Lemma \ref{lemma-faithful}.
Let $(V, \varphi)$ and $(W, \psi)$ be two descent data relative
to $X \to S$. Set $(V', \varphi') = f^*(V, \varphi)$ and
$(W', \psi') = f^*(W, \psi)$.
Let $\alpha' : V' \to W'$ be a morphism of descent data for $X'$ over $S$.
We have to show there exists a morphism $\alpha : V \to W$ of
descent data for $X$ over $S$ whose pullback is $\alpha'$.

\medskip\noindent
Recall that $V'$ is the base change of $V$ by $f$ and that
$\varphi'$ is the base change of $\varphi$ by $f \times f$
(see Remark \ref{remark-easier}).
By assumption the diagram
$$
\xymatrix{
V' \times_S X' \ar[r]_{\varphi'} \ar[d]_{\alpha' \times \text{id}} &
X' \times_S V' \ar[d]^{\text{id} \times \alpha'} \\
W' \times_S X' \ar[r]^{\psi'} &
X' \times_S W'
}
$$
commutes. We claim the two compositions
$$
\xymatrix{
V' \times_V V' \ar[r]^-{\text{pr}_i} &
V' \ar[r]^{\alpha'} &
W' \ar[r] &
W
}
, \quad i = 0, 1
$$
are the same. The reader is advised to prove this themselves rather
than read the rest of this paragraph. (Please email if you find a
nice clean argument.)
Let $v_0, v_1$ be points of $V'$ which map to the same point $v \in V$.
Let $x_i \in X'$ be the image of $v_i$, and let
$x$ be the point of $X$ which is the image of $v$ in $X$. In other words,
$v_i = (x_i, v)$ in $V' = X' \times_X V$. Write
$\varphi(v, x) = (x, v')$ for some point $v'$ of $V$.
This is possible because $\varphi$ is
a morphism over $X \times_S X$. Denote
$v_i' = (x_i, v')$ which is a point of $V'$.
Then a calculation (using the definition of $\varphi'$)
shows that $\varphi'(v_i, x_j) = (x_i, v'_j)$. Denote
$w_i = \alpha'(v_i)$ and $w'_i = \alpha'(v_i')$.
Now we may write $w_i = (x_i, u_i)$ for some point $u_i$ of $W$,
and $w_i' = (x_i, u'_i)$ for some point $u_i'$ of $W$.
The claim is equivalent to the assertion: $u_0 = u_1$.
A formal calculation using the definition of $\psi'$
(see Lemma \ref{lemma-pullback}) shows
that the commutativity of the diagram displayed above says that
$$
((x_i, x_j), \psi(u_i, x)) = ((x_i, x_j), (x, u'_j))
$$
as points of
$(X' \times_S X') \times_{X \times_S X} (X \times_S W)$
for all $i, j \in \{0, 1\}$. This shows that $\psi(u_0, x) = \psi(u_1, x)$
and hence $u_0 = u_1$ by taking $\psi^{-1}$.
This proves the claim because the argument above was formal
and we can take scheme points (in other words, we may
take $(v_0, v_1) = \text{id}_{V' \times_V V'}$).

\medskip\noindent
At this point we can use
Lemma \ref{lemma-fpqc-universal-effective-epimorphisms}.
Namely, $\{V' \to V\}$ is a fpqc covering as
the base change of the morphism $f : X' \to X$.
Hence, by
Lemma \ref{lemma-fpqc-universal-effective-epimorphisms}
the morphism $\alpha' : V' \to W' \to W$ factors through
a unique morphism $\alpha : V \to W$ whose base change is
necessarily $\alpha'$. Finally, we see the diagram
$$
\xymatrix{
V \times_S X \ar[r]_{\varphi} \ar[d]_{\alpha \times \text{id}} &
X \times_S V \ar[d]^{\text{id} \times \alpha} \\
W \times_S X \ar[r]^{\psi} & X \times_S W
}
$$
commutes because its base change to $X' \times_S X'$
commutes and the morphism $X' \times_S X' \to X \times_S X$
is surjective and flat (use Lemma \ref{lemma-ff-base-change-faithful}).
Hence $\alpha$ is a morphism of descent data
$(V, \varphi) \to (W, \psi)$ as desired.
\end{proof}

\end{document}
